% Options for packages loaded elsewhere
\PassOptionsToPackage{unicode}{hyperref}
\PassOptionsToPackage{hyphens}{url}
\PassOptionsToPackage{dvipsnames,svgnames,x11names}{xcolor}
\documentclass[
  10pt,
  fontset=fandol]{ctexart}
\usepackage{xcolor}
\usepackage[margin=16mm]{geometry}
\usepackage{amsmath,amssymb}
\setcounter{secnumdepth}{-\maxdimen} % remove section numbering
\usepackage{iftex}
\ifPDFTeX
  \usepackage[T1]{fontenc}
  \usepackage[utf8]{inputenc}
  \usepackage{textcomp} % provide euro and other symbols
\else % if luatex or xetex
  \usepackage{unicode-math} % this also loads fontspec
  \defaultfontfeatures{Scale=MatchLowercase}
  \defaultfontfeatures[\rmfamily]{Ligatures=TeX,Scale=1}
\fi
\usepackage{lmodern}
\ifPDFTeX\else
  % xetex/luatex font selection
\fi
% Use upquote if available, for straight quotes in verbatim environments
\IfFileExists{upquote.sty}{\usepackage{upquote}}{}
\IfFileExists{microtype.sty}{% use microtype if available
  \usepackage[]{microtype}
  \UseMicrotypeSet[protrusion]{basicmath} % disable protrusion for tt fonts
}{}
\makeatletter
\@ifundefined{KOMAClassName}{% if non-KOMA class
  \IfFileExists{parskip.sty}{%
    \usepackage{parskip}
  }{% else
    \setlength{\parindent}{0pt}
    \setlength{\parskip}{6pt plus 2pt minus 1pt}}
}{% if KOMA class
  \KOMAoptions{parskip=half}}
\makeatother
\setlength{\emergencystretch}{3em} % prevent overfull lines
\providecommand{\tightlist}{%
  \setlength{\itemsep}{0pt}\setlength{\parskip}{0pt}}
\usepackage{amsmath,amssymb,mathtools,needspace}
\xeCJKDeclareCharClass{CJK}{"2032,"0400->"04FF,"2160->"216B,"2460->"2473,"25A0,"25C7,"25CB,"25CF}
\IfFontExistsTF{Arial Unicode MS}{\xeCJKsetup{AutoFallBack=true}\setCJKfallbackfamilyfont{\CJKrmdefault}{Arial Unicode MS}}{}
\setcounter{secnumdepth}{0}
\setlength{\emergencystretch}{2em}
\allowdisplaybreaks
\usepackage{bookmark}
\IfFileExists{xurl.sty}{\usepackage{xurl}}{} % add URL line breaks if available
\urlstyle{same}
\hypersetup{
  pdftitle={最佳一致逼近多项式的存在性},
  colorlinks=true,
  linkcolor={Maroon},
  filecolor={Maroon},
  citecolor={Blue},
  urlcolor={Blue},
  pdfcreator={LaTeX via pandoc}}

\title{最佳一致逼近多项式的存在性}
\author{}
\date{}

\begin{document}
\maketitle

\subsubsection{3.2.1 最佳一致逼近多项式的存在性}\label{ux6700ux4f73ux4e00ux81f4ux903cux8fd1ux591aux9879ux5f0fux7684ux5b58ux5728ux6027}

Chebyshev 从另一观点研究一致逼近问题. 他不让多项式次数 \(n\) 趋于无穷,

\href{../../source-images/pdf-061.jpeg}{查看原页}

而是固定 \(n\), 记次数不大于 \(n\) 的多项式集合为 \(H_n\), 显然 \(H_n \subseteq C[a,b]\). 又记 \(H_n = \text{span}\{1,x,\cdots,x^n\}\), 其中 \(1,x,\cdots,x^n\) 是 \([a,b]\) 上一组线性无关的函数组, 是 \(H_n\) 中的一组基. \(H_n\) 中的元素 \(P_n(x)\) 可表示为 \[
P_n(x) = a_0 + a_1x + \cdots + a_nx^n,
\] 其中 \(a_0,a_1,\cdots,a_n\) 为任意实数. 要在 \(H_n\) 中求 \(P_n^*(x)\) 逼近 \(f(x) \in C[a,b]\), 使其误差 \[
\max_{a \leqslant x \leqslant b} |f(x) - P_n^*(x)| = \min_{P_n \in H_n} \max_{a \leqslant x \leqslant b} |f(x) - P_n(x)|,
\] 这就是通常所谓最佳一致逼近或 Chebyshev 逼近问题. 为了说明这一概念, 先给出以下定义.

定义 3.1 \(P_n(x) \in H_n, f(x) \in C[a,b]\), 称 \[
\Delta(f,P_n) = \|f - P_n\|_\infty = \max_{a \leqslant x \leqslant b} |f(x) - P_n(x)| \tag{3.2.1}
\] 为 \(f(x)\) 与 \(P_n(x)\) 在 \([a,b]\) 上的偏差.

显然, \(\Delta(f,P_n) \geqslant 0, \Delta(f,P_n)\) 的全体组成一个集合, 记作 \(\{\Delta(f,P_n)\}\), 它有下界 0. 若记集合的下确界为 \[
E_n = \inf_{P_n \in H_n} \{\Delta(f,P_n)\} = \inf_{P_n \in H_n} \max_{a \leqslant x \leqslant b} |f(x) - P_n(x)| \tag{3.2.2}
\] 则称 \(E_n\) 为 \(f(x)\) 在 \([a,b]\) 上的最小偏差.

定义 3.2 假定 \(f(x) \in C[a,b]\), 若存在 \[
P_n^*(x) \in H_n, \Delta(f,P_n^*) = E_n \tag{3.2.3}
\] 则称 \(P_n^*(x)\) 是 \(f(x)\) 在 \([a,b]\) 上的最佳一致逼近多项式或最小偏差逼近多项式, 简称最佳逼近多项式.

注意, 定义并未说明最佳逼近多项式是否存在, 但可证明下面的存在定理.

定理 3.2 若 \(f(x) \in C[a,b]\), 则总存在 \(P_n^*(x) \in H_n\), 使 \[
\|f(x) - P_n^*(x)\|_\infty = E_n.
\] 证明过程见文献{[}4{]}, 此处从略.

\end{document}
